\ifdefined\gkver\else\def\gkver{student}\fi
\documentclass[\gkver]{gaokaozhenti}
\begin{document}

\chapter{2007年天津理综}
%% paperType: 理综物理部分

\begin{enumerate}
\item
%% number: 14
%% paperName: 2007年天津理综第 14 题 6 分
%% typeId: 1
%% type: 单选题
%% chapter: 光学
%% point: 光的干涉
%% method: 无
%% score: 6
%% degree: 650
%% duplicateId: 0
%% body:
下列说法正确的是\xzanswer{B}

\fourchoices[answer=B]
{用三棱镜观察太阳光谱是利用光的干涉现象}
{在光导纤维束内传送图象是利用光的全反射现象}
{用标准平面检查光学平面的平整程度是利用光的偏振现象}
{电视机遥控器是利用发出紫外线脉冲信号来变换频道的}

%% answer: B
%% memo:
\memoanswer{
本题原卷及材料中未提供详解，待补充。
}

\item
%% number: 15
%% paperName: 2007年天津理综第 15 题 6 分
%% typeId: 1
%% type: 单选题
%% chapter: 运动和力
%% point: 动量和能量
%% method: 无
%% score: 6
%% degree: 650
%% duplicateId: 0
%% body:
如图所示，物体 A 静止在光滑的水平面上，A 的左边固定有轻质弹簧，与 A 质量相同的物体 B 以速度 $v$ 向 A 运动并与弹簧发生碰撞，A、B 始终沿同一直线运动，则 A、B 组成的系统动能损失最大的时刻是\xzanswer{D}

\onepicture[w=5.5cm]{figs/15.png}

\fourchoices[answer=D]
{A 开始运动时}
{A 的速度等于 $v$ 时}
{B 的速度等于零时}
{A 和 B 的速度相等时}

%% answer: D
%% memo:
\memoanswer{
本题原卷及材料中未提供详解，待补充。
}

\item
%% number: 16
%% paperName: 2007年天津理综第 16 题 6 分
%% typeId: 1
%% type: 单选题
%% chapter: 交流电
%% point: 交流电
%% method: 无
%% score: 6
%% degree: 650
%% duplicateId: 0
%% body:
将阻值为 $5\UO$ 的电阻接到内阻不计的交流电源上，电源电动势随时间变化的规律如图所示。下列说法正确的是\xzanswer{C}

\onepicture[w=5.0cm]{figs/16.png}

\fourchoices[answer=C]
{电路中交变电流的频率为 $0.25\UHz$}
{通过电阻的电流为 $\sqrt{2}\UA$}
{电阻消耗的电功率为 $2.5\UW$}
{用交流电压表测得电阻两端的电压是 $5\UV$}

%% answer: C
%% memo:
\memoanswer{
本题原卷及材料中未提供详解，待补充。
}

\item
%% number: 17
%% paperName: 2007年天津理综第 17 题 6 分
%% typeId: 1
%% type: 单选题
%% chapter: 曲线运动
%% point: 人造地球卫星
%% method: 无
%% score: 6
%% degree: 450
%% duplicateId: 0
%% body:
我国绕月探测工程的预先研究和工程实施已取得重要进展。设地球、月球的质量分别为 $m_{1}$、$m_{2}$，半径分别为 $R_{1}$、$R_{2}$，人造地球卫星的第一宇宙速度为 $v$，对应的环绕周期为 $T$，则环绕月球表面附近圆轨道飞行的探测器的速度和周期分别为\xzanswer{A}

\fourchoices[answer=A]
{$\sqrt{\frac{m_{2}R_{1}}{m_{1}R_{2}}}v$，$\sqrt{\frac{m_{1}R_{2}^{3}}{m_{2}R_{1}^{3}}}T$}
{$\sqrt{\frac{m_{1}R_{2}}{m_{2}R_{1}}}v$，$\sqrt{\frac{m_{2}R_{1}^{3}}{m_{1}R_{2}^{3}}}T$}
{$\sqrt{\frac{m_{2}R_{1}}{m_{1}R_{2}}}v$，$\sqrt{\frac{m_{2}R_{1}^{3}}{m_{1}R_{2}^{3}}}T$}
{$\sqrt{\frac{m_{1}R_{2}}{m_{2}R_{1}}}v$，$\sqrt{\frac{m_{1}R_{2}^{3}}{m_{2}R_{1}^{3}}}T$}

%% answer: A
%% memo:
\memoanswer{
本题原卷及材料中未提供详解，待补充。
}

\item
%% number: 18
%% paperName: 2007年天津理综第 18 题 6 分
%% typeId: 1
%% type: 单选题
%% chapter: 原子和原子核
%% point: 玻尔模型
%% method: 无
%% score: 6
%% degree: 650
%% duplicateId: 0
%% body:
右图为氢原子能级的示意图，现有大量的氢原子处于 $n=4$ 的激发态，当向低能级跃迁时辐射出若干不同频率的光。关于这些光下列说法正确的是\xzanswer{D}

\onepicture[w=3.6cm]{figs/18.png}

\fourchoices[answer=D]
{最容易表现出衍射现象的光是由 $n=4$ 能级跃到 $n=1$ 能级产生的}
{频率最小的光是由 $n=2$ 能级跃迁到 $n=1$ 能级产生的}
{这些氢原子总共可辐射出 $3$ 种不同频率的光}
{用 $n=2$ 能级跃迁到 $n=1$ 能级辐射出的光照射逸出功为 $6.34\UeV$ 的金属铂能发生光电效应}

%% answer: D
%% memo:
\memoanswer{
本题原卷及材料中未提供详解，待补充。
}

\item
%% number: 19
%% paperName: 2007年天津理综第 19 题 6 分
%% typeId: 1
%% type: 单选题
%% chapter: 磁场
%% point: 带电粒子在磁场中的运动
%% method: 无
%% score: 6
%% degree: 450
%% duplicateId: 0
%% body:
如图所示，在 $x$ 轴上方存在着垂直于纸面向里、磁感应强度为 $B$ 的匀强磁场。一个不计重力的带电粒子从坐标原点 O 处以速度 $v$ 进入磁场，粒子进入磁场时的速度方向垂直于磁场且与 $x$ 轴正方向成 $120^{\circ}$ 角，若粒子穿过 $y$ 轴正半轴后在磁场中到 $x$ 轴的最大距离为 $a$，则该粒子的比荷和所带电荷的正负是\xzanswer{C}

\onepicture[w=5.0cm]{figs/19.png}

\fourchoices[answer=C]
{$\frac{3v}{2aB}$，正电荷}
{$\frac{v}{2aB}$，正电荷}
{$\frac{3v}{2aB}$，负电荷}
{$\frac{v}{2aB}$，负电荷}

%% answer: C
%% memo:
\memoanswer{
本题原卷及材料中未提供详解，待补充。
}

\item
%% number: 20
%% paperName: 2007年天津理综第 20 题 6 分
%% typeId: 1
%% type: 单选题
%% chapter: 气体性质
%% point: 热力学第一定律
%% method: 无
%% score: 6
%% degree: 450
%% duplicateId: 0
%% body:
A、B 两装置，均由一支一端封闭，一端开口且带有玻璃泡的管状容器和水银槽组成，除玻璃泡在管上的位置不同外，其他条件都相同。将两管抽成真空后，开口向下竖直插入水银槽中（插入过程没有空气进入管内），水银柱上升至图示位置停止。假设这一过程水银与外界没有热交换，则下列说法正确的是\xzanswer{B}

\onepicture[w=3.2cm]{figs/20.png}

\fourchoices[answer=B]
{A 中水银的内能增量大于 B 中水银的内能增量}
{B 中水银的内能增量大于 A 中水银的内能增量}
{A 和 B 中水银体积保持不变，故内能增量相同}
{A 和 B 中水银温度始终相同，故内能增量相同}

%% answer: B
%% memo:
\memoanswer{
本题原卷及材料中未提供详解，待补充。
}

\item
%% number: 21
%% paperName: 2007年天津理综第 21 题 6 分
%% typeId: 1
%% type: 单选题
%% chapter: 机械波
%% point: 波的图象
%% method: 无
%% score: 6
%% degree: 450
%% duplicateId: 0
%% body:
如图所示，实线是沿 $x$ 轴传播的一列简谐横波在 $t=0$ 时刻的波形图，虚线是这列波在 $t=0.2\Us$ 时刻的波形图。已知该波的波速是 $0.8\Ums$，则下列说法正确的是\xzanswer{D}

\onepicture[w=6.0cm]{figs/21.png}

\fourchoices[answer=D]
{这列波的波长是 $14\Ucm$}
{这列波的周期是 $0.125\Us$}
{这列波可能是沿 $x$ 轴正方向传播的}
{$t=0$ 时，$x=4\Ucm$ 处的质点速度沿 $y$ 轴负方向}

%% answer: D
%% memo:
\memoanswer{
本题原卷及材料中未提供详解，待补充。
}

\item
%% number: 22
%% paperName: 2007年天津理综第 22 题 4 分
%% typeId: 4
%% type: 实验题
%% chapter: 电场
%% point: 示波器
%% method: 无
%% score: 4
%% degree: 450
%% duplicateId: 0
%% body:
\begin{enumerate}
	\item
	一种游标卡尺，它的游标尺上有 50 个小的等分刻度，总长度为 $49\Umm$。用它测量某物体长度，卡尺示数如图所示，则该物体的长度是\tkanswer[2.0cm]{4.120}$\Ucm$。

	\onepicture[w=8.5cm]{figs/22a.png}

	\item
	某学生用打点计时器研究小车的匀变速直线运动。他将打点计时器接到频率为 $50\UHz$ 的交流电源上，实验时得到一条纸带。他在纸带上便于测量的地方选取第一个计时点，在这点下标明 A，第六个点下标明 B，第十一个点下标明 C，第十六个点下标明 D，第二十一个点下标明 E。测量时发现 B 点已模糊不清，于是他测得 AC 长为 $14.56\Ucm$、CD 长为 $11.15\Ucm$，DE 长为 $13.73\Ucm$，则打 C 点时小车的瞬时速度大小为\tkanswer[2.0cm]{0.986}$\Ums$，小车运动的加速大小为\tkanswer[2.0cm]{2.58}$\Umsq$，AB 的距离应为\tkanswer[1.8cm]{5.99}$\Ucm$。（保留三位有效数字）

	\onepicture[w=10.0cm]{figs/22b.png}

	\item
	在“练习使用示波器”实验中，某同学将衰减调节旋钮置于最右边的“$\infty$”挡，扫描范围旋钮置于“外 X”档，“X 输入”与“地”之间未接信号输入电压，他在示波器荧光屏上看到的图象可能是下图中的\tkanswer[1.0cm]{B}。

	\fourpicture[numA=A, numB=B, numC=C, numD=D, labelprefix={}, w=2.6cm]
	{figs/22c.png}
	{figs/22d.png}
	{figs/22e.png}
	{figs/22f.png}
\end{enumerate}

%% answer: 
\jdanswer{
\begin{enumerate}
	\item
	$4.120\Ucm$

	\item
	$0.986\Ums$，$2.58\Umsq$，$5.99\Ucm$

	\item
	B
\end{enumerate}
}
%% memo:
\memoanswer{
本题原卷及材料中未提供详解，待补充。
}

\item
%% number: 23
%% paperName: 2007年天津理综第 23 题 16 分
%% typeId: 6
%% type: 计算题
%% chapter: 曲线运动
%% point: 竖直平面圆周运动
%% method: 无
%% score: 16
%% degree: 650
%% duplicateId: 0
%% body:
如图所示，水平光滑地面上停放着一辆小车，左侧靠在竖直墙壁上，小车的四分之一圆弧轨道 AB 是光滑的，在最低点 B 与水平轨道 BC 相切，BC 的长度是圆弧半径的 10 倍，整个轨道处于同一竖直平面内。可视为质点的物块从 A 点正上方某处无初速下落，恰好落入小车圆弧轨道滑动，然后沿水平轨道滑行至轨道末端 C 处恰好没有滑出。已知物块到达圆弧轨道最低点 B 时对轨道的压力是物块重力的 9 倍，小车的质量是物块的 3 倍，不考虑空气阻力和物块落入圆弧轨道时的能量损失。求：
\begin{enumerate}
	\item
	物块开始下落的位置距水平轨道 BC 的竖直高度是圆弧半径的几倍；
	\item
	物块与水平轨道 BC 间的动摩擦因数 $\mu$。
\end{enumerate}

\onepicture[w=5.5cm, align=right]{figs/23.png}

%% answer: 
\jdanswer{
\begin{enumerate}
	\item
	$h=4R$

	\item
	$\mu=0.3$
\end{enumerate}
}
%% memo:
\memoanswer{
（1）设物块开始下落处的位置距 $BC$ 的竖直高度为 $h$，圆弧轨道半径为 $R$。由机械能守恒定律得 $mgh=\frac{1}{2}mv^{2}$，在 $B$ 点根据牛顿第二定律得 $9mg-mg=m\frac{v^{2}}{R}$，解得 $h=4R$。

（2）物块滑到 $C$ 点时与小车的共同速度为 $v_{1}$，由动量守恒定律得 $mv=(m+3m)v_{1}$，对物块和小车应用动能定理得 $\mu mg\cdot 10R=\frac{1}{2}mv^{2}-\frac{1}{2}(m+3m)v_{1}^{2}$，解得 $\mu=0.3$。
}

\item
%% number: 24
%% paperName: 2007年天津理综第 24 题 18 分
%% typeId: 6
%% type: 计算题
%% chapter: 电磁感应
%% point: 电磁感应电路分析
%% method: 无
%% score: 18
%% degree: 450
%% duplicateId: 0
%% body:
两根光滑的长直金属导轨 MN、$M'N'$ 平行置于同一水平面内，导轨间距为 $l$，电阻不计，M、$M'$ 处接有如图所示的电路，电路中各电阻的阻值均为 $R$，电容器的电容为 $C$。长度也为 $l$、阻值同为 $R$ 的金属棒 ab 垂直于导轨放置，导轨处于磁感应强度为 $B$、方向竖直向下的匀强磁场中。ab 在外力作用下向右匀速运动且与导轨保持良好接触，在 ab 运动距离为 $s$ 的过程中，整个回路中产生的焦耳热为 $Q$。求：
\begin{enumerate}
	\item
	ab 运动速度 $v$ 的大小；
	\item
	电容器所带的电荷量 $q$。
\end{enumerate}

\onepicture[w=7.0cm, align=right]{figs/24.png}

%% answer: 
\jdanswer{
\begin{enumerate}
	\item
	$v=\frac{4QR}{B^{2}l^{2}s}$

	\item
	$q=\frac{CQR}{Bls}$
\end{enumerate}
}
%% memo:
\memoanswer{
（1）设 ab 上产生的感应电动势为 $E$，回路中电流为 $I$，ab 运动距离 $s$ 所用的时间为 $t$，则有 $E=Blv$，$I=\frac{E}{4R}$，$t=\frac{s}{v}$，$Q=I^{2}(4R)t$。由上述方程得 $v=\frac{4QR}{B^{2}l^{2}s}$。

（2）设电容器两极板间的电势差为 $U$，则有 $U=IR$，电容器所带电荷量为 $q=CU$，解得 $q=\frac{CQR}{Bls}$。
}

\item
%% number: 25
%% paperName: 2007年天津理综第 25 题 22 分
%% typeId: 6
%% type: 计算题
%% chapter: 运动和力
%% point: 动量定理
%% method: 无
%% score: 22
%% degree: 250
%% duplicateId: 0
%% body:
离子推进器是新一代航天动力装置，可用于卫星姿态控制和轨道修正。推进剂从图中 P 处注入，在 A 处电离出正离子，BC 之间加有恒定电压，正离子进入 B 时的速度忽略不计，经加速后形成电流为 $I$ 的离子束后喷出。已知推进器获得的推力为 $F$，单位时间内喷出的离子质量为 $J$。为研究问题方便，假定离子推进器在太空中飞行时不受其他外力，忽略推进器运动速度。
\begin{enumerate}
	\item
	求加在 BC 间的电压 $U$；
	\item
	为使离子推进器正常运行，必须在出口 D 处向正离子束注入电子，试解释其原因。
\end{enumerate}

\onepicture[w=7.0cm, align=right]{figs/25.png}

%% answer: 
\jdanswer{
\begin{enumerate}
	\item
	$U=\frac{F^{2}}{2JI}$

	\item
	推进器持续喷出正离子束，会使带有负电荷的电子留在其中，由于库仑力作用，将严重阻碍正离子的继续喷出，电子积累足够多时，甚至会将喷出的正离子再吸引回来，致使推进器无法正常工作。因此，必须在出口 D 处发射电子注入到正离子束，以中和正离子，使推进器获得持续推力。
\end{enumerate}
}
%% memo:
\memoanswer{
（1）设一个正离子的质量为 $m$、电荷量为 $q$，加速后的速度为 $v$，根据动能定理有 $qU=\frac{1}{2}mv^{2}$。设离子推进器在 $\Delta t$ 时间内喷出质量为 $\Delta M$ 的正离子，并以其为研究对象，推进器对 $\Delta M$ 的作用力为 $F'$，由动量定理有 $F'\Delta t=\Delta M v$；由牛顿第三定律知 $F'=F$。设加速后离子束的横截面积为 $S$，单位体积内的离子数为 $n$，则有 $I=nqvS$，$J=nmvS$，两式相比可得 $\frac{I}{J}=\frac{q}{m}$；又 $J=\frac{\Delta M}{\Delta t}$，解得 $U=\frac{F^{2}}{2JI}$。

（2）推进器持续喷出正离子束，会使带有负电荷的电子留在其中，由于库仑力作用，将严重阻碍正离子的继续喷出，电子积累足够多时，甚至会将喷出的正离子再吸引回来，致使推进器无法正常工作。因此，必须在出口 D 处发射电子注入到正离子束，以中和正离子，使推进器获得持续推力。
}

\end{enumerate}

\end{document}
